% =====================================================================
%  Dissertation front matter -- Kanak Chandra Dahal
%  Order: Abstract, Introduction, Literature Review, Methodology,
%  Objectives & Hypotheses, References (dissertation_refs.bib, all
%  entries verified against Crossref / arXiv / publisher records on
%  19 Nov 2026).
%  Compile: pdflatex -> bibtex -> pdflatex -> pdflatex
% =====================================================================
\documentclass[12pt,a4paper]{article}

\usepackage[margin=1.1in]{geometry}
\usepackage{amsmath,amssymb,amsthm}
\usepackage{times}

\usepackage{graphicx}
\usepackage{natbib}

\setlength{\parskip}{0.45em}
\setlength{\parindent}{0pt}

\title{\Large\bfseries Reconstructing Volatility Dynamics with Next
Generation Reservoir Computing:\\ A Simulation-Based Certification Study}
\author{Kanak Chandra Dahal\\
\small Atma Ram Sanatan Dharma College, University of Delhi, New Delhi}
\date{}

\begin{document}
\maketitle

% =====================================================================
\begin{abstract}
\noindent Financial volatility is persistent, asymmetric and partly driven
by the whole history of prices, yet most learned models of it are judged
only by how small their forecast error is. A model can forecast well and
still fail to represent the dynamics it has been fitted to. This
dissertation treats that gap directly. It asks whether Next Generation
Reservoir Computing (NG-RC), a data-efficient nonlinear model built from
polynomial features of delayed inputs, can reconstruct the effective
dynamics of volatility rather than merely predict its next value. NG-RC
is the first test case; the certification protocol itself is
model-agnostic and applies to any learned model of the dynamics. The
approach is a certification study. On synthetic data from processes whose
parameters are known, the trained model is tested against locked
acceptance criteria fixed before the data are generated, so that success
and failure are decided by the protocol and not adjusted afterwards. The
learned dynamics are then examined through their Jacobian: relaxation
times, stability and memory structure are read off the model itself and
compared with the known truth. Only after this controlled stage does the
work turn to realised volatility from market data, benchmarked against
HAR-RV, GJR-GARCH and standard machine-learning models. The aim is a modest but useful one: to find out
what a simple nonlinear model can and cannot recover about volatility
dynamics, with the failures measured as carefully as the successes.
\end{abstract}

\newpage

% =====================================================================
\section{Introduction}

Volatility is the quantity that matters for risk, and it is also the
quantity we cannot observe directly. Prices are recorded; volatility has
to be estimated, modelled, or inferred from options. Over sixty years of
empirical work have produced a stable list of things any serious model of
volatility must reproduce: returns are essentially unpredictable in sign,
but their magnitude is predictable; large moves cluster; the memory of
shocks decays very slowly; and losses move future volatility more than
gains of the same size \citep{mandelbrot1963,cont2001,ding1993,
bekaertwu2000}. The econometric tradition answers these facts with
parametric models, from ARCH and GARCH to the HAR family
\citep{engle1982,bollerslev1986,corsi2009}. The machine-learning tradition
answers them with flexible function approximators that often forecast
better \citep{bucci2020,christensen2023}.

Both traditions, however, share a habit that this dissertation questions:
models are evaluated almost entirely by forecast error. Forecast error is
a weak test of understanding. A model can achieve a low one-step error
while having completely wrong dynamics, and the difference becomes visible
only when the model is run autonomously, fed with its own output, or asked
to reproduce a quantity it was not directly trained on. For physical
systems this distinction is standard. Nobody would accept a climate model
that fits last week's temperature but produces the wrong seasonal cycle.
For volatility the distinction is rarely enforced, partly because the true dynamics are unknown and no clean experiment seems possible.

What ``the dynamics'' means for a stochastic process is worth one
sentence of care. The object is the conditional rule: given where the
system is now, the distribution of its next step is fixed. A GARCH
recursion is such a rule, and so is the path-dependent functional of
Section~2. Run forward, the rule generates every property of the
process, its persistence, its memory, its asymmetry. Reconstructing the
dynamics means recovering that rule from data, not merely fitting the
next value; and the rule is read off the trained model through exactly
the objects a physicist would use, the Jacobian and its spectra.

One can push back: if a model already forecasts with minimal error, why test anything else? Because one-step forecasts of daily realised volatility are nearly free. The series is so strongly autocorrelated that even a na\"ive forecast scores well, and forecast losses are too noisy to separate models reliably, which is why the benchmarking literature is full of close races and statistical ties \citep{grigoryeva2014,li2026,kilic2025}. The uses that matter depend on the dynamics themselves. Multi-day value-at-risk and stress tests compound the model over the horizon, so wrong persistence gives the wrong term structure of risk. Option prices depend on the distribution of the whole volatility path, not on point forecasts. And regime monitoring needs a notion of stability, which is a property of the dynamics, not of a forecast.

The way around this is to build the experiment artificially. If we
simulate a process whose dynamics are known exactly, fit a model to the
simulated data, and then ask what the model has learned, the gap between
learned and true dynamics can be measured rather than argued about. This
is the certification idea at the centre of the dissertation. The test is
fixed in advance: the acceptance criterion, the tolerance and the number
of replications are locked before any test data are generated, and the
model is not retuned afterwards. A model either passes the locked test or
it does not. At $n = 1{,}500$ replications the choice of confidence
bound is not cosmetic: a point estimate near the boundary can look
acceptable under a normal approximation while the exact binomial lower
bound still falls short. Under a locked protocol the exact bound is the
only defensible choice: the criterion is fixed in advance, and the
uncertainty of the pass rate itself is accounted for without
approximation. This design converts a vague question, whether a learned
model captures the dynamics, into a controlled experiment with a
pre-registered answer scale.

What should the locked test measure? A forecast error would only repeat
the usual mistake. The test should instead target an invariant: a number
fixed by the data-generating process, independent of any single
realisation. For a variance process the most basic invariant is
persistence, the rate at which a shock to volatility decays. In the
GARCH control it is exactly $\alpha+\beta$, and it fixes the relaxation
time over which the process forgets a shock. A learned model implies its
own persistence through the spectrum of its Jacobian, so the comparison
between learned and true dynamics reduces, in the first instance, to a
comparison of two numbers that are both known precisely. Later
invariants, the asymmetry of the response, the decay of memory, the
scaling of increments, follow the same pattern: each is a property of
the process, stated before the experiment, checked after it.

The model under examination is Next Generation Reservoir Computing
\citep{gauthier2021}. Reservoir computing in its classical form trains
only a linear readout on the states of a fixed random recurrent network
\citep{jaegerhaas2004,lukosevicius2009}; NG-RC removes the random network
altogether and constructs the state directly from delayed inputs and their
polynomial combinations. The result is a model that is nonlinear in the
data but linear in its parameters, trained by ridge regression in
closed form in seconds, with trained coefficients that can be read and
interpreted. This makes it
an unusually transparent candidate for reconstruction: the learned map is
explicit, and its Jacobian, the object through which relaxation times and
stability are computed, is available in closed form. NG-RC has matched
much larger models on chaotic benchmarks \citep{pathak2018,vlachas2020},
and reservoir methods have reached financial time series
\citep{grigoryeva2014,dossantos2025,li2026}, but the question asked here
has not been asked of it: not how well it forecasts volatility, but
whether what it learns is, in any defensible sense, the dynamics of
volatility.

Two questions organise the work. First, on controlled systems: does the
learned model, run on its own, reproduce the physical signatures of
volatility, meaning its relaxation timescales, its memory and scaling, and
its asymmetric response to gains and losses? Second, on the learned
dynamics themselves: does their stability, measured through the Jacobian,
carry information about changes in market regime at daily frequency?
The first
question is answered entirely on synthetic data, where truth is known.
The second moves to market data, and is meaningful only to the
extent that the first has established what the method can recover.

The dissertation is therefore best read as a measurement study. Its
contribution is not a new forecasting model but a way of testing learned
models against known dynamics, together with a concrete map of where one
simple, interpretable nonlinear model succeeds and where it fails.
Three things come out of the design: a certification protocol that any
learned model of the dynamics can be subjected to; a measurement of what
NG-RC in particular recovers and where it fails; and, when it fails, a
repair experiment in which candidate fixes are varied one at a time and
re-certified under the same locked criteria. The
literature review that follows collects the facts the model must
reproduce, the econometric and machine-learning benchmarks it must be
compared with, and the tools from nonlinear time series analysis through
which the comparison is made. The methodology section then builds the
experiment itself: the simulated ground truth, the locked acceptance
test, the diagnostics and the repair design. The final section states
the objectives and hypotheses precisely.

\newpage

% =====================================================================
\section{Literature Review}

\subsection{Stylised facts of volatility}

It is impossible to predict the direction in which prices will change, but
their magnitude can be predicted. As early as 1963, Mandelbrot noticed
that price moves were characterized by heavy-tailed distributions and
clustered large movements, and \citet{cont2001} later collected this and
other regularities into the standard list of stylised facts. For liquid
assets, with $r_t = \log P_t - \log P_{t-1}$, the autocorrelation
$\mathrm{Corr}(r_t, r_{t-k})$ is essentially zero at any lag $k \neq 0$,
while $\mathrm{Corr}(|r_t|, |r_{t-k}|)$ stays positive for a long time;
squared returns behave the same way. Yesterday's direction says nothing
about today's direction, but yesterday's magnitude says something about
today's variance. This is volatility clustering. The raw returns are
also heavy-tailed: their kurtosis far exceeds the Gaussian value, most
strongly at short horizons, and the excess shrinks under temporal
aggregation, exactly what a mixture of Gaussian shocks with
time-varying variance would produce \citep{cont2001}. The tails are
close to power laws, with tail indices near three across many markets,
a regularity stable enough to have earned a name, the inverse cubic
law \citep{cont2001}.

\citet{ding1993} examined powers of absolute returns and showed that the
autocorrelation of $|r_t|^d$ decays extremely slowly, particularly at
$d=1$. This property is called long memory, though its precise meaning
varies with the process and the estimator used. The statistical question
goes back to \citet{hurst1951}, whose rescaled-range analysis of reservoir
storage introduced the Hurst exponent, and to
\citet{grangerjoyeux1980}, who gave long memory a parametric form through
fractional differencing. \citet{lo1991} later warned how easily spurious
long memory can be read into stock prices, a warning any measurement of
memory has to respect. Formally, the slow decay is often written as a power law,
$\mathrm{Corr}(|r_t|, |r_{t-k}|) \sim k^{-\gamma}$ with a small exponent
$\gamma$, in contrast to the exponential decay $p^{k}$ that a
finite-state linear model produces. Either way, a model must create
persistence over long horizons without giving equal weight to all old
shocks. Figure~\ref{fig:garch} shows clustering in a controlled example:
a simulated GARCH$(1,1)$ path, in which calm and turbulent periods
alternate even though the mechanism is a single equation.

\begin{figure}[!htb]
\centering
\includegraphics[width=0.9\textwidth]{figs/fig_garch.pdf}
\caption{A simulated GARCH$(1,1)$ path ($\omega = 4\times10^{-6}$,
$\alpha = 0.06$, $\beta = 0.93$). Top: daily returns $r_t = \sigma_t
z_t$. Bottom: the conditional standard deviation $\sigma_t$. Large moves
cluster in time even though the shocks $z_t$ are independent; the
clustering is generated entirely by the variance recursion.}
\label{fig:garch}
\end{figure}

The second regularity is the leverage effect: losses increase future
volatility more than profits of equal size \citep{bekaertwu2000,
bouchaud2001}. The typical argument is that a drop in a company's stock
increases its leverage and makes the shares riskier, but the mechanism is
more complicated, and \citet{bekaertwu2000} argue for an alternative
volatility-feedback explanation. The mechanism does not matter here; what
matters is that the regularity exists. \citet{bouchaud2001} further found
that the reaction is lagged, so that past returns affect volatility
through time, which one can write schematically as
\begin{equation}
v_t = c + \sum_{j \geq 1} g_j\,(r_{t-j}),
\end{equation}
with $v_t$ a volatility parameter, $g_j$ the impact of the return $j$
periods ago, and $c$ a constant. This is a generalized distributed lag,
not an equation from either paper. Two further regularities complete the working list. Volatility itself
is asymmetric in its own moves: it rises faster after losses than it
falls after gains, so the volatility of volatility is not symmetric
either. And under temporal aggregation the whole structure relaxes
towards Gaussianity, so whatever mechanism one proposes must also
explain why the facts fade at lower frequencies \citep{cont2001}. Taken
together, clustering, long
memory and leverage say that volatility is determined by the entire price
history, not just the most recent return.

\subsection{The early econometric models}

The first econometric answer to clustering was ARCH \citep{engle1982}.
Writing $r_t = \sigma_t z_t$ with $z_t$ zero-mean and unit-variance, the
conditional variance is built from lagged squared returns,
\begin{equation}
\sigma_t^2 = \omega + \sum_{i=1}^{q} \alpha_i\, r_{t-i}^2 .
\end{equation}
\citet{bollerslev1986} added lagged variance to get GARCH:
\begin{equation}
\sigma_t^2 = \omega + \alpha r_{t-1}^2 + \beta \sigma_{t-1}^2,
\end{equation}
where $\omega > 0$ is constant, $\alpha$ the reaction to the most recent
squared return and $\beta$ the reaction to past variance. Iterating the recursion shows what the parameters mean. The
unconditional variance is $\bar\sigma^2 = \omega/(1-\alpha-\beta)$,
finite only when $\alpha+\beta<1$; in that case the process is
covariance-stationary and the effect of a single shock decays
geometrically, like $(\alpha+\beta)^k$. The sum $\alpha+\beta$ is
therefore the persistence of the process, and the half-life of a shock,
$\ln(1/2)/\ln(\alpha+\beta)$, grows without bound as the sum approaches
one. Typical daily estimates sit between $0.95$ and $0.995$, so a shock
takes weeks to months to halve. At the boundary $\alpha+\beta=1$ the
model becomes integrated GARCH, in which shocks never die, and the data
live close enough to that boundary that very persistent and permanent
are hard to tell apart with finite samples. Since nothing enters except squared
returns, plain GARCH is indifferent to the sign of the return. This was
fixed by Nelson's E-GARCH \citep{nelson1991} and by
\citet{gjr1993}, whose form is
\begin{equation}
\sigma_t^2 = \omega + \alpha r_{t-1}^2
+ \gamma\, \mathbf{1}_{\{r_{t-1} < 0\}}\, r_{t-1}^2
+ \beta \sigma_{t-1}^2 ,
\end{equation}
so that $\gamma > 0$ adds an extra effect of negative returns on the
conditional variance, the leverage effect that simple GARCH misses.
\citet{engleng1993} summarised this asymmetry through the news impact
curve, the map from yesterday's return to today's variance, and showed
that it is steeper for negative returns than for positive ones
(Figure~\ref{fig:news}). Nelson's EGARCH \citep{nelson1991} builds the
same effect in logarithms,
\begin{equation}
\log \sigma_t^2 = \omega + \beta \log \sigma_{t-1}^2
+ \alpha \left(|z_{t-1}| - \mathbb{E}|z|\right) + \gamma z_{t-1},
\end{equation}
which keeps the variance positive by construction and lets the sign of
the shock enter through $\gamma$. All of these models are estimated the
same way, by maximising a Gaussian quasi-likelihood; the estimates stay
consistent and asymptotically normal even when the true innovations are
not Gaussian, which is what makes the whole family practical
\citep{bollerslevwooldridge1992}.

\begin{figure}[!htb]
\centering
\includegraphics[width=0.62\textwidth]{figs/fig_news.pdf}
\caption{News impact curves for two variance recursions with the same
persistence. Plain GARCH responds only to the size of the shock; the
GJR form responds more strongly to negative returns, the leverage
asymmetry observed in equity data \citep{gjr1993,engleng1993}.}
\label{fig:news}
\end{figure}
Long memory entered the same family through the fractionally integrated
FIGARCH model of \citet{baillie1996}, in which the geometric decay of
the GARCH recursion is replaced by the hyperbolic weights of a
fractionally differenced shock, $(1-L)^{d}$ with $0 < d < 1$. \citet{englepatton2001} give a
sober assessment of what these models are and are not good for, and
\citet{taylor2007} remains a standard reference for the wider
econometrics of financial time series. These models showed that
volatility is a dynamic process rather than a sum of recent shocks, but
their functional form is fixed in advance.

\subsection{Realised volatility and the HAR family}

Daily closing prices only record the difference between two points in
time. \citet{andersen2003} recommended constructing the variance from
intraday returns instead. If $r_{t,j}$ is the $j$-th intraday return on
day $t$ and $M$ the number of intraday returns,
\begin{equation}
RV_t = \sum_{j=1}^{M} r_{t,j}^2,
\qquad
BV_t = \frac{\pi}{2}\,\frac{M}{M-1}\sum_{j=2}^{M} |r_{t,j}|\,|r_{t,j-1}| .
\end{equation}
Realised volatility is usually taken as $\sqrt{RV_t}$, though parts of the
literature call the variance itself ``realised volatility.'' Unlike the
squared daily return, $RV_t$ sees the variation inside the day even when
it cancels out by the close. Formally, as the sampling frequency grows, $RV_t$
converges to the integrated variance of the day, the time integral of
$\sigma^2(u)$ over trading hours, so realised volatility is a
measurement of the latent process rather than a model of it.
\citet{bn2004} tied realised variance to the
quadratic variation of the price process and introduced bipower variation,
which estimates the integrated continuous variance even in the presence of
finite-activity jumps. The approach had been prepared by
\citet{andersen1998} and \citet{andersen2001}, who answered the early
scepticism about GARCH-type models by showing that they forecast well
once volatility is measured properly, and
by \citet{hansenlunde2006}, who mapped out how microstructure noise
corrupts the intraday measurement and how to limit the damage. Their
lesson survives in the working practice of the field: sample sparsely
enough that the noise does not dominate, and treat the realised measure
itself as an estimate with error, an honesty that later forecast
comparisons must inherit \citep{patton2011}.

On top of this measurement, \citet{corsi2009} built the heterogeneous
autoregressive model HAR-RV, motivated by market participants acting on
different horizons:
\begin{equation}
RV_{t+1} = \beta_0 + \beta_d RV_t + \beta_w RV_t^{(w)}
+ \beta_m RV_t^{(m)} + \varepsilon_{t+1},
\end{equation}
\begin{equation}
RV_t^{(w)} = \frac{1}{5}\sum_{j=0}^{4} RV_{t-j},
\qquad
RV_t^{(m)} = \frac{1}{22}\sum_{j=0}^{21} RV_{t-j}.
\end{equation}
The model is linear with only a handful of parameters, yet it reproduces
most of the persistence of volatility, which is why it is the benchmark
of the field. The three terms form a crude cascade: each horizon's
coefficient loads on a different timescale, and the fitted sum
$\beta_d+\beta_w+\beta_m$ is typically close to one. A sum of a few
exponential decays can imitate a power law over the observed range of
lags, which is how HAR mimics long memory with four parameters. It has
no true long memory; it mimics long memory through
three horizons (Figure~\ref{fig:har}). Later work keeps the structure and changes the
information.

\begin{figure}[!htb]
\centering
\includegraphics[width=0.8\textwidth]{figs/fig_har.pdf}
\caption{The HAR-RV structure of \citet{corsi2009}. Tomorrow's realised
variance is a linear combination of yesterday's value, last week's
average and last month's average. Three horizons, four parameters, and
most of the persistence of the data.}
\label{fig:har}
\end{figure} \citet{pattonsheppard2015} split realised variation by the
sign of returns and showed that the components carry different information
about future volatility. \citet{harq2016} introduced the HARQ model, in
which the weight on daily variation depends on its own measurement error.
\citet{shephardsheppard2010} proposed the HEAVY family, which lets
realised measures follow their own GARCH-type recursion rather than
serving only as regressors, so that the measurement and its dynamics are
modelled together.

\subsection{Path-dependent and rough volatility}

\citet{guyonlekeufack2023} take a different route in ``Volatility is
(mostly) path-dependent.'' The history of index returns explains about
90\% of the variability of equity-index implied volatility and about 65\%
of the variability of noisy estimates of future realised daily volatility.
These numbers measure path dependence, not predictability of returns.
Their model is linear in two path functionals, one signed and one
built from magnitudes:
\begin{equation}
\sigma_t = \beta_0 + \beta_1 R_{1,t} + \beta_2 \sqrt{R_{2,t}},
\qquad
R_{1,t} = \sum_{j \geq 1} K^{(1)}_j\, r_{t-j}, \qquad
R_{2,t} = \sum_{j \geq 1} K^{(2)}_j\, r_{t-j}^2 ,
\end{equation}
with kernels given by shifted power laws over a range of memory lengths,
$\beta_1 < 0$ for leverage and $\beta_2 > 0$ for clustering. The
magnitude of past moves enters through the square root of a
kernel-weighted sum of squared returns, so the model is linear in a
volatility scale rather than in squared returns themselves. Volatility
is thus mostly a functional of how the market arrived at its present
state. The kernels $K^{(1)}$ and $K^{(2)}$ decay as shifted power
laws, so influence decays slowly but not uniformly: recent days matter
more, yet months-old moves still enter.

\citet{gatheral2018} arrive at a matching conclusion from a different
direction. Measuring the scaling of log-volatility increments, they find
$\mathbb{E}\,|\log\sigma_{t+\Delta}-\log\sigma_t|^{q} \propto
\Delta^{qH}$ with $H$ near $0.1$, far below the $H = 1/2$ of Brownian
motion. Modelling log-volatility as a fractional Ornstein--Uhlenbeck
process, $dX_t = -\kappa X_t\,dt + \nu\,dW_t^{H}$, reproduces the
scaling, and the pricing side of the same idea is developed by
\citet{bayer2016}, whose rough Heston model drives the variance with
fractional noise and derives its characteristic function from a
fractional Riccati equation, keeping the model tractable for option
pricing. Fractional noise with $H > 1/2$ is also the classical
route to long memory in stochastic volatility
\citep{comterenault1998}; roughness and long memory are the two ends of
the same dial. Small $H$ and power-law kernels are two descriptions of
one fact: volatility forgets its past neither quickly nor uniformly, and
the memory kernel matters more than any finite list of lags.
Figure~\ref{fig:kernel} makes the point concretely: over a year of daily
lags, a power-law kernel still assigns visible weight to months-old
moves where a geometric, GARCH-like kernel has already forgotten them.
Figure~\ref{fig:rough} shows the scaling measurement itself on synthetic
paths whose roughness is known: the fitted slopes recover the true
exponents, which is what makes the empirical finding $H \approx 0.1$
believable.

\begin{figure}[!htb]
\centering
\includegraphics[width=0.7\textwidth]{figs/fig_kernel.pdf}
\caption{Normalised weights on past returns for two power-law kernels
and a geometric kernel with persistence $0.95$, on a log--log scale.
Geometric decay is indistinguishable from a power law over the first few
weeks and then falls off far faster; long-memory behaviour is precisely
what lives in the tail of this plot.}
\label{fig:kernel}
\end{figure}

\begin{figure}[!htb]
\centering
\includegraphics[width=0.62\textwidth]{figs/fig_rough.pdf}
\caption{The signature plot behind rough volatility, computed here on
two simulated paths whose Hurst exponents are known. The mean absolute
increment grows as a power of the lag, and the fitted slopes recover
$H = 0.1$ and $H = 1/2$ respectively. On market log-volatility the
same measurement gives slopes near $0.1$ \citep{gatheral2018}.}
\label{fig:rough}
\end{figure}

\subsection{Option-implied volatility}

Options give a second, forward-looking observation of volatility. The
pricing framework of \citet{blackscholes1973} and \citet{merton1973}
turns an option price into an implied volatility by inversion, and
\citet{heston1993} gave the standard stochastic-volatility correction
with its closed-form characteristic function. Implied volatilities vary
across strikes and maturities in a surface with its own stable shape
\citep{gatheral2006}, and the SVI parametrisation of
\citet{gatheraljacquier2014} gives each maturity slice a parsimonious
form with checkable arbitrage conditions. In its raw form the total
implied variance at log-moneyness $k$ is
\begin{equation}
w(k) = a + b\left(\rho\,(k-m) + \sqrt{(k-m)^2 + s^2}\right),
\end{equation}
five parameters per maturity slice, and the absence of butterfly and
calendar arbitrage translates into explicit inequalities on them. The implied surface lives under
the risk-neutral measure, realised volatility under the real-world one,
and the difference between the two is the volatility risk premium: the
compensation investors demand for bearing volatility risk, which makes
implied volatility sit systematically above realised volatility on
average. The premium is not noise; it is a priced, persistent object
with its own dynamics. Checking a reconstruction claim on both sides is
the natural extension of the present framework; this dissertation works
entirely on the real-world side, and the risk-neutral side is left to
future work.

\subsection{Machine learning for volatility}

Machine learning entered volatility modelling mainly through forecasting.
\citet{bucci2020} applied neural networks to realised volatility, and
\citet{christensen2023} compared regularised regressions, trees and
neural networks against HAR models. Flexible models could beat HAR even
when fed only lagged daily, weekly and monthly realised variances, with
larger advantages at longer horizons, where nonlinear mappings represent
persistence better. Their variable-importance analysis is just as
instructive: most of the predictable structure sits in the same daily,
weekly and monthly aggregates that HAR already uses, so the nonlinear
models add shape to the response rather than new information. Recurrent architectures such as the LSTM
\citep{lstm1997} are the natural deep-learning candidates for such data:
their gated state is designed to carry information over long horizons,
which is exactly what volatility demands. HARNet \citep{harnet2022} and
DeepVol \citep{deepvol2024} build the multiple time scales directly into
the architecture, using dilated convolutions whose receptive field spans
daily to monthly lags, a neural echo of the HAR cascade. The
evidence is not one-sided: \citet{kilic2025}, in a Federal Reserve study
comparing HAR, ARFIMA and regime-switching HAR variants with gradient
boosting and several neural architectures on S\&P~500 and equity-level
realised volatility, find that machine learning sometimes improves on HAR
but does not systematically beat the broader econometric set, and that
rankings depend on the asset, horizon, estimation window and loss
function. Where the machines do win, they win by little; where they
lose, the loss is systematic. Flexibility alone is not an argument. \citet{horvath2021} use
networks differently, as fast proxies for calibrating rough volatility
models, without changing the underlying mathematics. Machine learning can
improve forecasting, accelerate existing models, or find structure; this
dissertation is concerned with the third use.

\subsection{Reservoir computing and NG-RC}

Reservoir computing is another option. In the echo state networks of
Jaeger and the liquid state machines of \citet{maass2002}, a fixed,
randomly wired recurrent network maps the input series into
high-dimensional states and only the output weights are trained
\citep{jaegerhaas2004,lukosevicius2009,nakajima2021}. In equations, an echo state network keeps a state
$\mathbf{x}(t) = f\!\left(W\,\mathbf{x}(t-1) + W_{\mathrm{in}}\,
\mathbf{u}(t)\right)$ with fixed random $W$ and $W_{\mathrm{in}}$, and
the trained object is only $\hat{\mathbf{y}}(t) = W_{\mathrm{out}}\,
\mathbf{x}(t)$. The echo state property, guaranteed in practice by
keeping the spectral radius of $W$ below one, says that the state is
asymptotically a function of the input history alone
\citep{jaegerhaas2004}. The price is a
large set of hyperparameters (reservoir size, spectral radius, sparsity,
input scaling), dependence on the random initialisation, and a warm-up
period. \citet{gauthier2021} removed the random network and replaced it
with a nonlinear vector autoregression built from polynomial features of
delayed inputs, calling the construction Next Generation Reservoir
Computing. On the standard chaotic benchmarks they report forecasts
competitive with classical reservoirs orders of magnitude larger, at a
small fraction of the training cost, which is what makes the
construction attractive for a replication-heavy study. Given $\mathbf{X}_t \in \mathbb{R}^d$, the delay vector with
$k$ delays spaced by $s$ is
\begin{equation}
\mathbf{O}_{\mathrm{lin}}(t) = \mathbf{X}_t \oplus \mathbf{X}_{t-s} \oplus
\cdots \oplus \mathbf{X}_{t-(k-1)s},
\end{equation}
which has $dk$ elements ($\oplus$ is concatenation). The quadratic library
holds the unique pairwise products, $dk(dk+1)/2$ of them rather than
$(dk)^2$ since $x_i x_j = x_j x_i$, and the prediction is a linear
readout:
\begin{equation}
\mathbf{O}_{\mathrm{total}}(t) =
\begin{bmatrix}
1 \\ \mathbf{O}_{\mathrm{lin}}(t) \\ \mathbf{O}_{\mathrm{nonlin}}^{(2)}(t)
\end{bmatrix},
\qquad
\hat{\mathbf{X}}_{t+1} = \mathbf{W}_{\mathrm{out}}\,
\mathbf{O}_{\mathrm{total}}(t).
\end{equation}
The model is linear in its parameters but nonlinear in the input, and
$\mathbf{W}_{\mathrm{out}}$ is computed by ridge regression to balance fit
against stability. Writing $O$ for the matrix whose columns are the
feature vectors and $Y$ for the targets, the ridge solution is
$W_{\mathrm{out}} = Y\,O^{\top}\left(O\,O^{\top} +
\beta I\right)^{-1}$, a single linear solve. Figure~\ref{fig:ngrc}
summarises the architecture and the free-running mode in which the
prediction is fed back as the next input.

\begin{figure}[!htb]
\centering
\includegraphics[width=0.95\textwidth]{figs/fig_ngrc.pdf}
\caption{The NG-RC architecture. The observed series is packed into a
delay vector, expanded into constant, linear and quadratic features, and
mapped to the next value by a linear readout trained with ridge
regression. In free-running mode the model's own prediction re-enters as
input, which is where wrong dynamics become visible.}
\label{fig:ngrc}
\end{figure} Compared with a classical reservoir there is far less
noise and no warm-up, and the trained coefficients can be analysed
directly. On chaotic benchmarks the record is strong: reservoir
computers reproduce the spatiotemporal dynamics of the
Kuramoto--Sivashinsky equation far beyond the training horizon
\citep{pathak2018}, infer unmeasured variables of a chaotic system from
the measured ones \citep{lu2017}, and remain competitive when the
readout is trained by backpropagation through time instead
\citep{vlachas2020}. Physical implementations are
an active field of their own \citep{tanaka2019}.

Reservoirs have reached financial data, but the treatment is incomplete.
\citet{grigoryeva2014} used time-delay reservoirs to forecast conditional
covariances of GARCH-type synthetic processes and realised volatilities
computed from intraday quotes, with forecast error as the measure of
success. Their finding was nuanced: the reservoir was competitive with
econometric benchmarks but not decisively better, and the ranking moved
with the dataset and the error measure, an early sign that one-step
error is a coarse instrument for these questions. \citet{dossantos2025} combine a classical reservoir with the
Hurst exponent, computed by rescaled-range analysis over overlapping
100-day windows, to forecast next-day closing values of eleven stock
indices, reporting smaller errors and larger $R^2$ than LSTM, ARIMA,
bidirectional LSTM and transformer models in most cases. The study
connects reservoirs with long-range dependence, but the target is the
price level rather than volatility, $R^2$ on such a persistent series is
not informative by itself, no na\"ive forecast is included in the
comparison, and the hyperparameter choices differ across models.
\citet{li2026} push in yet another direction, forecasting realised
volatility with a quantum reservoir; notably, their quantum model is
statistically tied with classical benchmarks on the standard confidence
sets, so even there the interesting question is what the model has
learned rather than its headline error. A tie in forecast error between
two models with completely different internal dynamics is the cleanest
possible evidence that the error alone cannot tell us what a model
knows. All of this work ends at forecast
error: what does the reservoir know, and how does it behave when left to
run on its own?

\subsection{Tools from nonlinear time series analysis}

Asking what a model has learned requires tools that compare dynamics, not
just errors. Nonlinear time series analysis supplies them
\citep{kantzschreiber2004}. Takens' embedding theorem
\citep{takens1981} shows that the geometry of a deterministic system can
be reconstructed from a single observed series through delay vectors:
for a scalar $x_t$, the points $(x_t, x_{t-\tau}, \dots,
x_{t-(m-1)\tau})$ unfold the attractor up to a smooth change of
coordinates once $m$ is large enough, and
\citet{grassberger1983} gave the standard algorithm for measuring the
dimension of the reconstructed attractor. Lyapunov exponents measure the
rate at which nearby states separate: the largest,
$\lambda_{\max} = \lim_{t \to \infty} \frac{1}{t} \ln |\delta\mathbf{x}_t
/ \delta\mathbf{x}_0|$, is positive for chaos and negative for
contracting dynamics, and finite-time exponents give a local,
time-resolved version of the same quantity along a trajectory. For a
map, each exponent is a long-run average of the stretching produced by
the product of Jacobians along an orbit, so the Jacobian of a learned
map is not an arbitrary object to study: it is exactly the object
through which dynamical-systems theory reads stability
\citep{kantzschreiber2004}. The spectral radius of one Jacobian is the
local, instantaneous analogue of the same measurement, which is the form
this dissertation uses. Memory can be quantified
through autocorrelation decay, the Hurst exponent \citep{hurst1951}, or
detrended fluctuation analysis \citep{peng1994}, which measures the
growth of fluctuations $F(n) \propto n^{a}$ of the integrated series
over windows of length $n$, with $a > 1/2$ signalling persistent
correlations. The correlation dimension, in particular,
estimates how the number of attractor points inside a ball scales with
the ball's radius, and a saturating value as the embedding dimension
grows is the classical signature of low-dimensional dynamics
\citep{grassberger1983,kantzschreiber2004}. These quantities are
properties of the process, not of a particular forecast, which is exactly
what a certification test needs. Agent-based work such as
\citet{luxmarchesi1999} also reminds us that clustering and heavy tails
can emerge from simple interacting components, so reproducing stylised
facts is necessary but not sufficient evidence that a model has the right
mechanism. A final tool guards against reading structure into noise:
the surrogate-data method of \citet{theiler1992} generates artificial
series that share the data's linear properties but destroy any
nonlinear structure, and a measured quantity is trusted only if it
differs between the data and its surrogates. The same logic, compare
against a null that preserves the trivial explanation, runs through the
whole methodology of this dissertation.

\subsection{How forecasts are compared}

Even the forecasting part of this dissertation needs better accounting
than a single error table. \citet{dieboldmariano1995} give the standard
test for equal predictive accuracy between two models: for the loss
series $d_t = L(e_t^{(1)}) - L(e_t^{(2)})$ formed from the two models'
errors, the statistic $\bar d / \sqrt{\widehat{\mathrm{Var}}(\bar d)}$,
with a long-run variance estimate, is asymptotically standard normal
under the null;
\citet{hansen2005} extends the comparison to many models: his test for
superior predictive ability asks whether any model in a pool beats a
benchmark, and bootstraps studentised loss differences so that the data
mining a large comparison invites is priced into the null. The model
confidence set of \citet{hansenlundenason2011} takes the final step:
rather than crowning one winner, it eliminates models until only a set
of statistically indistinguishable ones remains, which is the honest
answer whenever several models tie. And
\citet{patton2011} shows which loss functions remain consistent when
volatility itself is only observed through an imperfect proxy such as
realised variance: with an unbiased proxy, squared error and the QLIKE
loss, $\log \hat v + \tilde\sigma^2 / \hat v$ for a variance forecast
$\hat v$ and proxy $\tilde\sigma^2$, rank competing forecasts the same
way the unobservable truth would. When the object of interest is a risk
measure rather than a point forecast, loss comparison gives way to
coverage testing: the binomial exceedance test of \citet{kupiec1995},
the standard backtest for value-at-risk, is a close statistical cousin
of the exact binomial bound this dissertation certifies with. These
results set the ground rules for any claim that
one model forecasts better than another, and they apply with extra force
to machine-learning comparisons, where the number of tried architectures
is usually large and unreported.

\subsection{The gap}

The issue is not that the field needs another forecast. A model can
forecast excellently and still fail to reproduce the process behind the
persistence or asymmetry of volatility, and a good one-step forecast says
little about what happens when the model's own output is fed back as
input. Even the studies that combine reservoirs with nonlinear dynamics,
such as \citet{dossantos2025}, measure quality in pricing errors.

What the literature suggests instead is a combination of techniques in
which no single technique is treated as a complete description of
volatility. ARCH and GARCH give simple controls; HAR-RV gives a reliable
multi-horizon benchmark; path-dependent models give history-based
nonlinear feedback with signed returns; NG-RC provides a mechanism by
which such terms become part of an explicit polynomial map. It is open
whether the shared mathematical structure of NG-RC and path-dependent
models is enough to reproduce the known properties from noisy data. The
purpose here is neither to prove that markets are generated by a
low-dimensional deterministic system nor to claim trading profits. The
question is whether a simple, data-efficient nonlinear model can identify
the structure of volatility dynamics, simulate it on its own, and provide
a meaningful way to measure stability. The plan is to start with
simulations of a path-dependent process, with GARCH as a control, to test
whether parameters can be recovered; then to move to empirical volatility,
with daily, weekly and monthly lags including consecutive ones, benchmarked
against HAR-RV and other econometric and machine-learning models; and
finally to study the trained parameters, the behaviour of the simulated
series, and the usefulness of a stability index derived from Jacobian
analysis. Framed this way, the study sits inside the classical
system-identification tradition \citep{ljung1999}: a dynamical system
is learned from its output, and the learned model is then required to
reproduce the system's key properties. The difference is that here the
model class itself is learned rather than parametric, so validation
cannot be inherited from the class and has to be designed.

Simulation is what makes this program possible. On real data, evaluation
confounds two errors: the model's error and the measurement error of the
volatility proxy, since realised measures are themselves noisy estimates
of a latent object \citep{andersen1998,andersen2001}. A model can look
right because the proxy is wrong. In simulation the true $\sigma_t$ is
known exactly, so the comparison isolates estimation error, the only
error a certification study is about. Table~\ref{tab:map} collects the
map this review has drawn.

\begin{table}[!htb]
\centering
\small
\begin{tabular}{p{2.6cm}p{4.0cm}p{3.6cm}p{3.4cm}}
\hline
Approach & What it assumes & What it explains well & What it cannot
show \\
\hline
GARCH family & one-step recursion in past news and variance & clustering,
persistence, leverage (with sign terms) & long memory; whether a learned
model has the same dynamics \\
HAR-RV & three fixed horizons of
realised variance & most of the persistence with four parameters & any
mechanism; it is a regression, not a generator \\
Rough / path-
dependent & power-law memory kernels, fractional noise & the scaling of
volatility increments; joint SPX--VIX smiles & a trained model's internal
stability \\
Deep learning & flexible recurrence or dilated
convolutions & small, inconsistent forecast gains & what the network
has learned about the process \\
NG-RC & polynomial features of
delays, linear readout & interpretable map; closed-form Jacobian &
untested: the open question of this dissertation \\
\hline
\end{tabular}
\caption{The state of the field in one view. Each row summarises a
family reviewed above; the last column is the gap this work addresses.}
\label{tab:map}
\end{table}

\newpage

% =====================================================================
\section{Methodology}

The work has three stages. First, the model is certified on synthetic
data where the truth is known. Second, any failure is taken apart with
diagnostics that separate bias from noise. Third, the model is repaired,
one factor at a time, and re-certified under the same protocol.
Figure~\ref{fig:pipeline} shows the full loop; this section describes
each stage in turn. Table~\ref{tab:notation} collects the notation.

\begin{table}[!htb]
\centering
\small
\begin{tabular}{cl}
\hline
Symbol & Meaning \\
\hline
$\sigma_t^2$ & conditional variance of the return at time $t$ \\
$\varepsilon_t$ & return innovation, $\varepsilon_t = \sigma_t z_t$ \\
$z_t$ & standardised innovation, mean zero, unit variance \\
$\omega,\alpha,\beta$ & GARCH(1,1) level, news and persistence parameters \\
$p$ & GARCH persistence $\alpha+\beta$, the certified quantity \\
$k$ & number of delays in the NG-RC feature vector (44 in the validation run) \\
$\rho$ & target spectral radius of the delay block (the NG-RC $\gamma$) \\
$\lambda$ & ridge penalty in the linear readout \\
$J$ & Jacobian of the trained map at a trajectory point \\
$\rho(J)$ & spectral radius of $J$; the certified persistence \\
$q_L$ & one-sided lower 95\% bound on the per-setting recovery rate \\
\hline
\end{tabular}
\caption{Notation. The first block belongs to the data-generating
process; the rest belong to the estimator and the certificate.}
\label{tab:notation}
\end{table}

\begin{figure}[!htb]
\centering
\includegraphics[width=0.98\textwidth]{figs/fig_pipeline.pdf}
\caption{The certification loop. Ground truth is simulated, the model is
trained under frozen settings, and the learned Jacobian is tested
against a locked acceptance criterion. A failure is diagnosed, repaired
one factor at a time, and re-certified under the same protocol; a pass
moves the protocol to the next model or the next invariant.}
\label{fig:pipeline}
\end{figure}

\subsection{Ground truth by simulation}

The controlled stage uses two generators. The first is a GARCH$(1,1)$
control,
\begin{equation}
\varepsilon_t = \sigma_t z_t, \qquad
\sigma_t^2 = \omega + \alpha\,\varepsilon_{t-1}^2 + \beta\,\sigma_{t-1}^2,
\qquad z_t \sim \mathrm{i.i.d.},
\end{equation}
whose persistence $p = \alpha + \beta$ is known exactly. Its
unconditional variance is $\bar\sigma^2 = \omega/(1-\alpha-\beta)$ and
its relaxation time is $\tau = -1/\ln(\alpha+\beta)$, so the parameter
the test targets also controls how quickly shocks die. The second is the
path-dependent volatility process of \citet{guyonlekeufack2023}, in
which the volatility is a direct functional of the whole return path.
In the discrete form used for simulation this reads
\begin{equation}
r_t = \sigma_t z_t, \qquad
\sigma_t = \beta_0 + \beta_1 \sum_{j \geq 1} K^{(1)}_j\, r_{t-j}
+ \beta_2 \sqrt{\sum_{j \geq 1} K^{(2)}_j\, r_{t-j}^2},
\qquad z_t \sim \mathrm{i.i.d.},
\end{equation}
with power-law kernels: the history enters directly,
not through a one-step recursion, and the randomness enters through the
return path itself. Both are choices,
not conveniences. GARCH is the benchmark the econometrics literature
already understands, and the path-dependent process is the nonlinear
alternative whose feedback structure NG-RC should, in principle, be able
to express. The four frozen settings of the validation run are all
GARCH, with persistence values near $0.95$ and $0.98$, close to what is
estimated on daily equity data; the Guyon--Lekeufack cell follows the
identical protocol. One difference is stated plainly: the GARCH
recursion carries its persistence in closed form, the path-dependent
generator does not, so the recovery target for that cell is derived
analytically, by linearising the generator around its stationary point.
Write $S_t = \sum_{j \geq 1} K^{(2)}_j\, x_{t-j}$ for the magnitude
channel, with $x_t = r_t^2$, $\bar S$ its stationary mean and
$\sigma^* = \beta_0 + \beta_2 \sqrt{\bar S}$ the stationary volatility
level. Linearising the square root about $\bar S$ and conditioning on
the squared-return path alone, the signed channel drops out at first
order because the innovations are symmetric, and the one-step
conditional mean of $x$ becomes a linear functional of the delay vector,
\begin{equation}
\mathbb{E}[x_{t+1} \mid x_t, x_{t-1}, \dots] \;\approx\;
\mathrm{const} + \sum_{j \geq 1} d_j\, x_{t-j+1},
\qquad
d_j = \frac{\sigma^*\,\beta_2}{\sqrt{\bar S}}\, K^{(2)}_j
+ \beta_1^2\,\bigl(K^{(1)}_j\bigr)^2 .
\end{equation}
The two terms are the two channels of the generator: the magnitude
kernel enters linearly, the signed kernel through its square, because
only magnitudes are visible in $x$. The stationary moments are pinned by
the parameters through a scalar fixed-point equation,
$\bar x\,\bigl(1 - \beta_1^2 \sum_j (K^{(1)}_j)^2\bigr) =
\bigl(\beta_0 + \beta_2 \sqrt{\bar x \sum_j K^{(2)}_j}\bigr)^2$,
so every $d_j$ is a function of the generator's parameters alone and no
simulation enters the definition. The recovery target is the dominant
relaxation rate of this linear map: the largest root $p^*$ of
\begin{equation}
\sum_{j \geq 1} d_j\, p^{-j} = 1,
\end{equation}
the exact analogue of $\alpha + \beta$ for the control. It is computed
once from the frozen parameters and locked with the protocol before any
model is trained. Because the parameters of both processes are set by hand,
every recovery question has an exact answer, and no part of the ``truth''
needs to be estimated. Two working rules keep the simulation honest.
Each dataset is generated after a burn-in long enough for the process to
forget its initial condition, so that no estimate is contaminated by
start-up transients. And every random seed is fixed and recorded, so
that any figure or number in the study can be reproduced exactly.

The design of the controlled stage belongs to the simulation-based
calibration tradition \citep{talts2020sbc}: when ground truth is
available by construction, a procedure should be judged by whether it
recovers that truth under a fixed criterion, applied identically in
simulation and in deployment. The protocol of the next section is the
frequentist version of that idea. Volatility paths are simulated from a
known generator, the learned model must recover a known dynamical
quantity at a rate not meaningfully below a true-model positive
control, and the pass count is read against an exact binomial law, with
the tolerance calibrated to what the control achieves. \citet{talts2020sbc} warn that calibration
checks averaged over a prior can miss failures that appear only in
particular regions of the data. The protocol answers that warning
structurally, because nothing is averaged: each generator, each
persistence level and each invariant carries its own certificate, so a
conditional failure is a failed certificate, not a diluted average.

\subsection{A locked certification protocol}

The central design choice is that the acceptance criteria are fixed
before any data are generated, and are not adjusted afterwards. A tolerance fixed in absolute terms hides a
sampling problem: near the unit root even the best estimator of
persistence is imprecise, so an absolute band can be too tight for the
truth itself to pass. The protocol therefore runs a positive control
alongside the model \citep{boger2026}. For the GARCH control this is
the quasi-maximum-likelihood fit of the true model, put through the
identical pipeline: the same series length, the same split and the same
number of replications. The locked tolerance is calibrated to what the
control achieves, its half-width set from the control's own sampling
spread, so that the control passes at a pre-decided reference rate.
Certification is then non-inferiority to that control: the model's
estimates must land inside the calibrated band in more than $90\%$ of
independent datasets in every pre-decided setting, the same requirement
the control itself meets. Success is recovery of the parameter: the
estimate must sit inside the band, and no forecast error of any kind
enters the criterion. The required rate is checked with an exact binomial
confidence bound \citep{clopperpearson1934}, not a normal approximation,
so that the uncertainty of the pass rate itself is accounted for. The
bound is one-sided by choice: the certificate fails in only one
direction, a recovery rate that is too low, so the error budget is
spent where the risk lies rather than split across two tails.
Concretely, if $s$ of $n$ independent datasets pass, the one-sided lower
$95\%$ bound $q_L$ is defined by
\begin{equation}
\sum_{j=s}^{n} \binom{n}{j}\, q_L^{\,j}\, (1-q_L)^{\,n-j} = 0.05,
\end{equation}
and certification requires $q_L > 0.90$ in every pre-decided setting.
The size of the run was chosen against this rule. At $n = 1{,}500$ the
bound clears $0.90$ once at least $1{,}370$ datasets fall inside the
band, a pass fraction of $91.3\%$, and the bound's half-width around a
$90\%$ rate is about $1.5$ percentage points, fine enough to separate a
true recovery rate of $90\%$ from one of $93\%$. The test is strict in
both directions: a model that truly recovers the parameter $89\%$ of the
time passes with probability below one in five hundred, while a model
that truly recovers $93\%$ of the time passes with probability $99.4\%$.
Because the band is calibrated so that the control's own recovery rate
sits at the reference level, clearing the same bound is a
non-inferiority statement: the model's recovery rate lies within the
bound's resolution of the best rate any estimator achieves. The earlier
absolute tolerance, ten percent of the gap $1-p$, is kept as a
descriptive diagnostic of how far the estimates sit from the truth in
absolute terms; it no longer gates certification.
Figure~\ref{fig:band} illustrates the test schematically. The
validation run used four frozen settings with $1{,}500$ independent
datasets each, $6{,}000$ in total. Each dataset is an independent
simulated series of $T = 20{,}000$ daily steps, generated after a
burn-in of $1{,}000$ steps; the first $70\%$ of a series trains the
readout and the remainder evaluates it. The two counts are easy to
confuse, so they are stated once, plainly: $1{,}500$ is the number of
replications per setting, not the length of any series. Nothing was
retuned between settings,
and no run was stopped early or repeated because of its result. An
earlier pilot was kept strictly exploratory: it informed the design, and
the setup was frozen before the fresh validation began. Read in
econometric terms, the certificate is a moment condition with an exact
sampling law under replication. Moment-based estimation of volatility
models is classical \citep{zhangwang2023}; the difference is that the
moment certified here is dynamical, the recovery of a known parameter,
rather than distributional.

\begin{figure}[!htb]
\centering
\includegraphics[width=0.95\textwidth]{figs/fig_band.pdf}
\caption{Schematic of the acceptance test at $p = 0.95$. Each point is
one replication's estimate $\hat{p}$; the shaded band is the locked
tolerance, calibrated to the positive control, and the drawn width is
illustrative rather than to scale. Left: the pattern a
certified model would show. Right: the pattern a failing model shows, estimates centred below
the truth and scattered far outside the band.}
\label{fig:band}
\end{figure} The point of
these rules is that success and failure are decided by the protocol, so
that a failure means something about the model rather than about the
analyst's patience.

\subsection{The estimator and the Jacobian}

The model under test is NG-RC \citep{gauthier2021}. The input is a fixed
vector of past values of one series, and in the controlled stage that
series is the squared returns $x_t = r_t^2$ the generator emits, the
realised measure of the simulation. The latent variance path defines
the target and is never shown to the model: the only noise in the input
is the process's own innovation, there is no further measurement error,
and the noise-contaminated case is taken up with the scope
limitations. The feature library collects
constant, linear and quadratic combinations of these delays; and the
readout is a single linear map trained by ridge regression, which has a
closed-form solution and one penalty hyperparameter. All hyperparameters,
the number of delays, the training length and the penalty, are chosen
once and frozen across every setting of the validation, so that any
difference in outcomes is a property of the data-generating process and
not of the tuning. Training is
therefore cheap and exactly reproducible, which is what makes large
replication studies like the one above practical. The validation run
reported here used $k = 44$ consecutive daily delays with linear
features. Writing the trained map as $\hat x_{t+1} = f(x_t,
\dots, x_{t-k+1})$, the Jacobian entries are $J_{ij} = \partial f
/ \partial x_{t-j+1}$; because $f$ is polynomial in its arguments,
each entry is available in closed form, and for a quadratic library each
entry is an affine function of the current state. The conditional
expectation $\mathbb{E}[x_{t+1} \,|\, x_t, x_{t-1}, \dots]$ is the
conditional variance path itself, so a delay map on squared returns
approximates the same contraction the GARCH recursion encodes in
$\alpha+\beta$; this is the sense in which the spectral radius
$\rho(\hat J)$ targets the persistence. The learned relaxation time is
$\hat\tau = -1/\ln\rho(\hat J)$, directly comparable with the
$\tau = -1/\ln(\alpha+\beta)$ of the control.

The certificate comes from the learned dynamics, not the forecast error.
The spectral radius is the right first object for this. A map that
forgets its initial condition must contract perturbations, which is
$\rho < 1$; the same condition, one ring further along, is the echo
state property that makes classical reservoirs well-defined
\citep{jaegerhaas2004}, and the approximation theory that licenses
reading dynamics off a contracting, input-driven map is by now rigorous
\citep{harthookdawes2021,hart2024}. And how far $\rho$ sits below one is the
persistence: it sets the rate at which shocks to the learned system die,
which is precisely the quantity $\alpha+\beta$ encodes in the control.
Because the trained map is an explicit polynomial in the delays, its
Jacobian can be computed in closed form at any point of a trajectory.
The persistence estimator is the spectral radius of this Jacobian,
\begin{equation}
\hat{p} \;=\; \rho\!\left(\widehat{J}\right),
\end{equation}
which plays the role that $p = \alpha + \beta$ plays in the control. A
model that has truly reconstructed the dynamics should return $\hat{p}$
near the true $p$; a model that has merely learned a good one-step
forecast need not.

Three clarifications keep the estimator honest. First, it never
free-runs: the Jacobian is evaluated at states of the observed
validation series, and no multi-step rollout enters the certificate, so
the well-known divergence of closed-loop simulations is irrelevant to
it. Second, the certificate compares the learned radius with the known
recursion parameter of a simulated process; it is not an estimate of a
Lyapunov exponent of a stochastic system, a quantity with its own strict
definition through products of Jacobians along a trajectory
\citep{nelson1990}. Third,
where the learned map is used to generate paths rather than to read a
parameter, modelling log-volatility keeps the variance positive by
construction. The certificate, however, belongs to the volatility map:
$\alpha+\beta$ is the contraction rate of the variance recursion, and a
map trained in log space has a different Jacobian, whose spectral
radius is not $\alpha+\beta$. The validation therefore trains and reads
the estimator on the squared-return series; a log-space variant would first
need its target redefined in log space, measured from the true
generator's log-space response.

Why eigenvalues at all is a fair question, and it has a principled
answer. A reservoir computer with fixed nonlinear features and a
trained linear readout is a dictionary of observables whose induced
dynamics are linear in the readout, which is exactly the structure of a
finite-dimensional approximation of the Koopman operator
\citep{gulinamauroy2021,guyangshi2026}. Under this reading the
eigenvalues of the learned map approximate Koopman eigenvalues, the
growth and decay modes of the dynamics, so reading relaxation times off
the Jacobian spectrum is not a heuristic but the natural spectral
object of the model class. The reservoir's memory, controlled by its
spectral radius, determines which of these modes remain observable
\citep{guyangshi2026}. The operator view is stated for deterministic
flows, and the volatility processes here are noise-driven; the claim is
only that the learned map plays the role of the effective deterministic
skeleton of the system, the framing used throughout this chapter.

The spectral certificate also has an econometric reading. In a
GARCH$(1,1)$ the impact of a variance shock decays at the rate
$\alpha+\beta$, and the volatility impulse response function
\citep{hafnerherwartz2006,hafnerherwartz2023} makes that decay the
object of study in parametric form. The dominant eigenvalue of the
learned Jacobian is the same quantity estimated nonparametrically: a
one-period volatility impulse response, read off the fitted map instead
of a GARCH likelihood. The certificate that $\rho(\hat{J})$ recovers
$\alpha+\beta$ within the locked band is therefore a data-driven VIRF
check, with the closed-form GARCH response as its parametric benchmark.
The asymmetric responses of
\citet{hafnerherwartz2023}, derived for multivariate GARCH and read
here at the univariate level, are the
impulse-response counterpart of the leverage asymmetry the remaining
invariants test.

\subsection{Diagnostics: separating bias from width}

When a model fails, the failure is split into two parts before any
repair is attempted. Given $n$ replications with estimates $\hat p_1,
\dots, \hat p_n$, write $\bar p$ for their mean and $s$ for their
sample standard deviation. The first part is a shifted centre: the mean
estimate sits away from the truth, measured in Monte Carlo standard
errors, $|\bar p - p| / (s/\sqrt{n})$, so that ``away'' has a
quantitative meaning. The second part is excess width: $s$ itself is
large compared with the half-width of the tolerance band.
Figure~\ref{fig:diag} shows the two failure shapes. A shifted centre points to a
systematic distortion, such as shrinkage; excess width points to missing
structure, such as truncated memory. A third possibility is irreducible
sampling noise: near the unit root even the best estimator scatters,
and the positive control measures that floor directly, so width is read
as missing structure only when it exceeds what the control achieves.
Outside the controlled stage a further channel appears: a noisy
realised-volatility proxy attenuates the estimated persistence, which
shifts the centre downward. Two robustness checks accompany the
decomposition: the tolerance band is halved and doubled to confirm the
conclusion does not hinge on the exact tolerance, and an independent
pilot is rerun to confirm the failure replicates. Persistence near one
is a classical hard inference problem, and the standard econometric
answers, median-unbiased estimation \citep{andrews1993} and the grid
bootstrap \citep{hansen1999grid}, are built for least-squares
autoregressions. The estimator here is a ridge-shrunk spectral radius
of a companion-like matrix, so those pivots do not transfer directly;
any bootstrap in this setting must resample the ridge estimator itself.

\begin{figure}[!htb]
\centering
\includegraphics[width=0.95\textwidth]{figs/fig_diagnostics.pdf}
\caption{The diagnostic split, schematic. Left: a shifted centre, where
the estimates cluster tightly but away from the truth. Right: excess
width, where the estimates are centred correctly but scatter far beyond
the locked band (shaded). The two failures call for different repairs.}
\label{fig:diag}
\end{figure} These diagnostics are
reported as post-hoc analysis. They were not part of the locked
criteria, and they are used to explain failure, not to reclassify it.

One caveat belongs with any eigenvalue-based readout. The Jacobian of a
delay-embedding model is companion-like and strongly non-normal, and in
non-normal systems perturbations can grow transiently even when every
eigenvalue lies inside the unit circle. Recent work shows this is not a
corner case: non-normal amplification in random multiplicative processes
can generate heavy-tailed, clustered fluctuations far from spectral
instability \citep{troude2025,troude2026}. The certificate is protected
at the control level, where the truth it compares against is a scalar
contraction rate, but the caveat bites when the estimator is read as a
description of transient behaviour. A trajectory-level growth check,
tracking how the learned map actually amplifies perturbations over
several steps, is the stricter diagnostic, and it is recorded here as a
planned addition rather than a claim.

\subsection{Repair as a designed experiment}
\label{sec:repair}

\begin{figure}[!htb]
\centering
\includegraphics[width=0.95\textwidth]{figs/fig_repair.pdf}
\caption{The repair experiment. Each candidate repair is applied on its
own, and the repaired model re-enters the same locked protocol. A pass
measures the repair; a fail sharpens the negative result.}
\label{fig:repair}
\end{figure}

Repairs are tested one factor at a time. If several things change at
once, an improvement cannot be attributed, and a failure cannot be
diagnosed further. Three repairs follow from the diagnostics, and each
carries its own testable signature: a repair is believed only if the
failure pattern changes in the direction its hypothesis predicts.
Table~\ref{tab:repairs} summarises the design.

\begin{table}[!htb]
\centering
\small
\begin{tabular}{p{2.9cm}p{2.6cm}p{4.1cm}p{4.2cm}}
\hline
Repair & Targets & What changes & Signature if the hypothesis is right \\
\hline
1. Ridge sweep + debiasing & shifted centre & penalty on a log grid;
bias measured against known truth & centre moves to the truth as the
penalty falls; width unchanged \\
2. Memory scan & excess width & delay
span and layout; HAR-type aggregated inputs & scatter narrows; centre
unchanged \\
3. Objective change & centre and width & multi-step or
invariant-matching training loss & long-horizon behaviour improves at
similar one-step fit \\
\hline
\end{tabular}
\caption{The repair experiment at a glance. Each repair is applied
alone, aimed at one failure mode, and read through the change it should
produce if its hypothesis is correct.}
\label{tab:repairs}
\end{table}

\paragraph{1. Ridge shrinkage and debiasing.} The ridge penalty
contracts the weights, and through them the spectrum of the learned
Jacobian, so persistence is biased downward. The first repair sweeps the
penalty on a logarithmic grid and retrains at each value. Because the
truth is known in simulation, the bias can be measured directly as a
function of the penalty, giving an empirical bias curve from which a
debiased estimate is read off. If lowering the penalty moves the
estimated persistence towards the truth, shrinkage is confirmed as the
distortion; if it does not, the hypothesis is rejected cheaply and the
factor is ruled out. This is the simulation version of a
classical idea: the downward small-sample bias of estimated
autocorrelations, and its analytic corrections, have been known in the
time-series literature for seventy years
\citep{marriottpope1954,kendall1954}.

\paragraph{2. Memory structure.} A GARCH process remembers its whole
past through its recursion, while a finite delay vector remembers only
its own window. The second repair varies the delay span and compares
short, long and non-uniform delay layouts; if longer or better-placed
memory narrows the scatter of the estimates, truncation is confirmed as
the source of the excess width. The linear memory capacity
measure of \citet{jaeger2002} gives a way to quantify how much of the
past a finite reservoir can retain at all. Non-uniform layouts of the
daily--weekly--monthly kind used in Objective 2, or HAR-type aggregated
inputs \citep{corsi2009}, inject long memory at the input side without
growing the model.

\paragraph{3. Training objective.} One-step error rewards the next value,
not the long-run dynamics. The third repair replaces or augments the
one-step loss: multi-step losses penalise the error of free-running
rollouts, and invariant-matching penalties require simulated paths to
match the autocorrelation or the marginal distribution of the data.
Closely related stabilisation ideas are documented in the reservoir
literature on chaotic and climate systems, including hybrid and
correction schemes \citep{pathak2018,wikner2020}, systematic studies of
which design choices control free-running performance \citep{platt2022},
and noise-inspired regularisation that stabilises long-horizon
prediction \citep{wikner2024}.

Each repair is also a measurement. If debiasing alone moves the centre,
shrinkage was the story, and the fix is cheap. If the memory scan
narrows the scatter, the delay layout was too poor for the persistence
being tested, and the repair cost is paid in design, not in data. If
neither moves the numbers, the feature library itself lacks the
nonlinearity the dynamics need, which is the strongest and most
informative failure of the three, because it localises the limitation in
the model class rather than in its settings.

\subsection{Recertification}

Every repaired model re-enters the same locked protocol, with the same
settings, the same band and the same binomial bound. Improvement is not
a better-looking forecast; it is a pass rate that meets the pre-decided
requirement. Trying several repairs and reporting only the one that
passes would quietly reopen the data-mining problem the locked design
exists to close, so each repair is evaluated on its own, with its own
binomial bound, and all outcomes are reported. This closes the loop: test, diagnose, repair, re-certify.
If a repaired model passes, the protocol has measured the repair. If it
still fails, the failure is now sharper, because the diagnostics have
already ruled repairs in or out. Either outcome is a result. The loop
then generalises: the same locked test can be pointed at any other
learned model, an LSTM, a dilated-convolution network, a quantum
reservoir, and each model's certificate or failure is comparable with
every other's, because the yardstick never moved.

\subsection{Scope and limitations}

The protocol certifies recovery on the generators it is pointed at, and
no further. GARCH$(1,1)$ and the Guyon--Lekeufack process represent two
mechanisms, a one-step recursion and direct history dependence; they are
test beds with known truth, not claims about what markets are. A pass
says the model reconstructs these mechanisms; it says nothing about
mechanisms not tested.

Two further boundaries are stated plainly. The estimator and the
certificate concern conditional-mean dynamics and their persistence;
nothing here claims to recover the full conditional law of volatility.
Distributional questions, quantiles, predictive intervals and tails,
belong to future work with quantile or generative readouts, where the
natural certificates are predictive-calibration checks and proper
scoring rules such as the CRPS. And the
certification stage runs on simulated series where the realised measure
is observed exactly as generated, so no measurement-noise channel exists
inside the experiment;
extending the protocol to noise-contaminated observations, through
filtering or a noise-robust loss, is a second line of future work.

The tolerance is calibrated to the positive control rather than fixed
in absolute terms. An absolute band, ten percent of the gap $1-p$ in
the pilot design, bounds the distortion of the relaxation time, the
quantity a physicist would read off the dynamics, but it can sit below
the best estimator's own sampling spread near the unit root, in which
case even the true model fails for reasons that have nothing to do with
the model. Calibrating to the control removes that floor. The protocol
does not depend on the particular calibration; it depends on the
calibration being frozen. Likewise, the estimator tested here
is one point in the NG-RC design space, a linear library over $k=44$
consecutive delays with frozen hyperparameters. Conclusions about NG-RC
as a class require the repair loop and wider sweeps, and the protocol is
built to make those follow-ups cheap and comparable.

Certification is per-invariant. A model can recover persistence and fail
asymmetry, or pass on the GARCH control and fail on the path-dependent
generator, and each outcome is a separate certificate. Finally, the
real-data stages inherit the measurement error of realised measures:
Objective 2 onwards compares against proxies, not truth, which is
exactly why the controlled stage comes first and why its verdicts carry
the most weight.

Finally, a matter of process. The protocol, the locked criterion and
the code are archived with timestamps in the project repository, which
fixes their order relative to the results, but they were not entered in
an independent public registry. That is a gap in form rather than in
substance, and it is recorded here as one.

\newpage

% =====================================================================
\section{Objectives and Hypotheses}

\subsection{Overall aim}

This dissertation investigates whether NG-RC can serve as a data-driven
reconstruction of the effective deterministic dynamics of financial
volatility, treated as a noise-driven nonlinear dynamical system. NG-RC
is, however, only the first test case. The deeper object of study is the
certification protocol itself: a model-agnostic procedure that tests a
learned model of the dynamics on data where the truth is known, takes
its failures apart with diagnostics, repairs them as a designed
experiment, and re-certifies the result under unchanged criteria. What
is demonstrated here on NG-RC can in principle be repeated on any model
whose learned dynamics can be read off, through a Jacobian or an
equivalent description. Rather than judging the model by one-step
forecast error alone, it asks two questions, pursued respectively
through Objective 1 and Objectives 2 and 3 below.

\begin{enumerate}\renewcommand{\labelenumi}{\textbf{Q\arabic{enumi}.}}
\item Does the learned model reproduce the physical
signatures of volatility: relaxation timescales, memory and scaling, and
the asymmetric response to gains and losses?
\item Does the stability of the learned dynamics, measured through its
Jacobian, carry meaningful information about changes in market regime at
daily frequency?
\end{enumerate}

\subsection{Specific objectives}

\paragraph{Objective 1: Validate recovery of dynamics on controlled
systems.}
Generate synthetic series from a GARCH process, as a control, and from a
path-dependent volatility process of the Guyon--Lekeufack type, both
written in a common Langevin-type stochastic-dynamics form. Test whether
NG-RC recovers the known relaxation rates, feedback asymmetry and memory
structure, and whether Jacobian-based relaxation times and contraction
rates match the analytically known values. This step establishes what
the method can and cannot recover from the noisy squared-return series
of a known generator before it is applied to real markets.

\paragraph{Objective 2: Reconstruct empirical volatility dynamics with a
physically motivated delay structure.}
Apply NG-RC to realised volatility, constructed from intraday returns,
using daily, weekly and monthly delays, and consecutive delays, as a
non-uniform delay embedding. Benchmark against HAR-RV, at least one
econometric model (GJR-GARCH) and at least one machine-learning model,
evaluated by out-of-sample forecast accuracy and by the interpretability
of the trained coefficients.

\paragraph{Objective 3: Construct and evaluate a physically grounded
stability index.}
Define a stability index from the Jacobian of the trained model, using
the spectral radius and the local contraction rates read off the
Jacobian along the trajectory, and compute it on rolling windows. Assess whether the
index varies systematically across two stress windows fixed in
advance, the 2008 financial crisis and the March 2020 COVID shock, with
lead time and false-alarm rate as the pre-registered metrics, reported
against a rolling realised-volatility threshold as the primary
benchmark. The
index is treated as a descriptive diagnostic of the learned dynamics, not
as a trading or prediction tool.

\paragraph{Objective 4: Diagnose, repair and re-certify the failures
found in Objective 1.}
The arc of the dissertation does not end at the test. Where the
certification of Objective 1 fails, apply the diagnostics of Section 3.4
to separate a shifted centre from excess width, and then run the repair
experiment of Section 3.5: a ridge sweep with measured debiasing, a
memory scan over delay layouts, and changes to the training objective,
varied one factor at a time so that each effect is attributable.
Re-certify every repaired model under the same locked protocol, so that
improvement is defined by the pass rate and not by inspection. A repair
that works, and a repair that fails in an informative way, are both
results of the study.

\subsection{Hypotheses}

\begin{enumerate}\renewcommand{\labelenumi}{\textbf{H\arabic{enumi}.}}
\item (Objective 1) The spectral radius of the learned Jacobian recovers
the known persistence $\alpha + \beta$ of the GARCH control to within
estimation error.
\item (Objective 2) NG-RC with a signed quadratic library is
out-of-sample competitive with HAR-RV, and a library with odd terms
outperforms an even-only library in capturing the return--volatility
asymmetry.
\item (Objective 3) The stability index rises systematically in the
pre-registered stress windows, the 2008 crisis and the March 2020 shock,
with lead time and false-alarm rate fixed in advance and reported
against a rolling realised-volatility threshold as the primary
benchmark.
\item (Objective 4) Under the unchanged locked protocol, the diagnostics
attribute each certification failure to a specific cause, a shifted
centre or excess width, and at least one of the three repairs then moves
the GARCH-control pass rate monotonically upward, without relaxing the
criterion. Crossing the locked acceptance band is sufficient for success
but not required: an attributed, monotone improvement that still falls
short of the band is a result of the study.
\end{enumerate}

\subsection{Closing remark: what to do when a model fails}

A certification protocol is useful only if failure can be acted on. The
methodology of Section~\ref{sec:repair} treats repair as a designed
experiment: a ridge sweep with measured debiasing for a shifted centre,
a memory scan for excess width, and multi-step or invariant-matching
objectives when the training loss itself is the distortion. Each
repaired model is then re-certified under the same locked protocol, so
that improvement is measured rather
than asserted.

% =====================================================================
\newpage
\bibliographystyle{plainnat}
\bibliography{dissertation_refs}

\end{document}
